\section{Introduction}
\label{sec:introduction}

Asteroids preserve evidence of the formation and subsequent evolution of
the Solar System. Their orbital distribution reflects the gravitational
influence of the planets and the redistribution of material during the
early evolution of the planetary system, while their compositions and
physical properties record different formation environments and histories
of collisions and thermal processing. The main belt is particularly useful
for studying these processes because its dynamical structure can be
examined together with the properties of a large population of individual
objects \citep{DeMeoCarry2014,Morbidelli2015Belt}. Connecting these
observations to models of Solar System evolution requires both reliable
orbits and measurements that describe the diversity within the asteroid
population.

Optical astrometry and photometry provide complementary parts of this
information. Repeated positions establish and improve orbits, connect
observations obtained at different apparitions, and enable targeted
observations of individual objects \citep{Jedicke2015Surveys}. Brightness
variations constrain rotation periods and amplitudes; with sufficient
sampling of rotational phase and viewing geometry, photometry can also
support estimates of spin orientation and shape
\citep{WarnerHarrisPravec2009,Durech2015Models}. These applications place
different demands on temporal sampling. Measurements distributed across a
night trace apparent motion and short-term brightness changes, whereas
observations on additional nights and at different apparitions broaden
the rotational and geometric coverage. Survey photometry is therefore
valuable both for assembling time series of individual asteroids and for
selecting objects for more intensive follow-up.

Wide-field surveys have greatly increased the scale on which such
measurements can be obtained. Pan-STARRS combines repeated imaging with
attribution and linking of moving-object detections
\citep{Denneau2013}; ATLAS emphasizes frequent coverage of large sky areas
for nearby-object searches \citep{Tonry2018}; and the Zwicky Transient
Facility combines a large field of view with repeated observations for
time-domain astronomy \citep{Bellm2019}. Dedicated moving-object analysis
also makes use of the observing patterns of broader time-domain surveys,
as demonstrated by the Heliocentric-orbiting Objects Processing System for
the Wide Field Survey Telescope \citep{Wang2025HOPS}. Together, these
programs illustrate how instrument design, observing cadence, and
association methods shape the resulting small-body observations.

The scientific use of a survey depends on the coverage and sampling
achieved during observations. Limiting magnitude and image quality affect
which objects are measured, while sky position, solar elongation, and
apparent angular rate influence the populations represented in the
sample. The interval between repeated exposures determines the
displacement available for recognizing motion; the distribution of
measurements across subsequent nights determines the length and density
of the astrometric and photometric time series. These selections also
enter the interpretation of asteroid population statistics
\citep{Jedicke2015Surveys}. Describing a survey's footprint, cadence, and
measurement quality provides the basis for selecting appropriate
observations and for understanding which scientific questions its data
can address.

We developed \SHARP\ to carry out repeated ecliptic imaging with the
Xinglong 60/90~cm Schmidt telescope. The approximately
$2.94^\circ\times2.95^\circ$ field of view allows many asteroids to be
observed in each pointing. The standard observing sequence uses 30~s
white-light exposures and aims for three visits to a field during a night,
sampling angular motion over a baseline of roughly an hour. Repeated
visits on later nights build longer sequences for objects that remain
within the survey footprint. This strategy is intended to provide a
continuing set of asteroid positions and brightness measurements, with
particular emphasis on the main-belt population accessible to the
telescope.

The same observations support association with the predicted positions of
known asteroids and a search for coherent motion among source detections.
The related GOTTA Prototype study uses known-asteroid measurements to
characterize survey observations \citep[in preparation]{Xiao2026InPrep}.
Known objects supply positional references for evaluating astrometric
residuals and allow the number and temporal distribution of measurements
to be determined for each identity. They also provide a direct measure of
the fraction of repeatedly detected objects that are recovered as
single-night linkages. The motion search supplies candidate tracks for
image inspection and subsequent identification. In this paper, these
measurements are used to describe the performance of the survey and the
observational content of its first year, establishing the basis for
subsequent studies of asteroid variability and for follow-up of selected
objects.

We present the survey design and its first year of observations, from
2025 November~15 to 2026 November~15. Section~\ref{sec:survey} describes the telescope,
calibration, and observing strategy. Section~\ref{sec:processing}
introduces the moving-object analysis. Section~\ref{sec:performance}
examines the sky coverage, cadence, photometric and astrometric quality,
and linkage of repeated detections. Section~\ref{sec:samples} describes
the known-asteroid observations and single-night candidates.
Section~\ref{sec:discussion} discusses the scientific uses of these
observations and summarizes the first-year survey.
